\begin{figure}[htbp]
\centering
\begin{tikzpicture}[>=Stealth]
  \draw (-3.35,-1.55) rectangle (-0.15,1.40);
  \fill (-1.75,-.05) circle (1.5pt) node[left] {$0$};
  \draw[fill=blue!8] (-1.75,-.05) circle (0.82);
  \node at (-1.75,-1.30) {$B_\varepsilon(0)$};
  \draw[->] (-1.75,-.05) -- (-1.07,.53) node[above] {$tv$};
  \draw[->] (-1.75,-.05) -- (-2.40,.58) node[above] {$tw$};
  \node at (-1.75,-1.92) {$T_pM$};
  \draw[->] (.15,0) -- node[above] {$\exp_p$} (1.65,0);
  \draw (1.95,-1.4) .. controls (3.0,-1.85) and (5.4,-1.55) .. (5.9,-.35) .. controls (6.25,.65) and (5.0,1.55) .. (3.55,1.23) .. controls (2.55,1.02) and (1.8,.0) .. (1.95,-1.4);
  \fill[blue!8] (3.55,-.10) circle (0.90);
  \fill (3.55,-.10) circle (1.5pt) node[below] {$p$};
  \node at (2.45,-1.56) {$U$};
  \draw[->,thick] (3.55,-.10) .. controls (4.05,.18) and (4.30,.53) .. (4.55,.93);
  \draw[->,thick] (3.55,-.10) .. controls (3.05,.30) and (2.83,.62) .. (2.93,.96);
  \draw[densely dashed] (4.8,-.95) .. controls (5.35,-.52) and (5.43,.13) .. (4.92,.72);
  \node[align=center,font=\scriptsize] at (5.20,-1.55) {possible cut or\\conjugate behavior};
  \node at (4.10,1.62) {$M$};
\end{tikzpicture}
\caption{A normal neighborhood is the image of a small tangent ball under $\exp_p$, where radial geodesics give local coordinates. The dashed curve marks a possible global obstruction: outside the shaded domains $\exp_p$ need not be injective and radial geodesics need not remain minimizing.}
\label{fig:vii23-c15-exponential-scope}
\end{figure}

\begin{table}[htbp]
\centering\small
\begin{tabular}{p{0.27\linewidth}p{0.61\linewidth}}
\toprule
Property & What it asserts (and what it does not assert)\\
\midrule
Normal neighborhood & $\exp_p$ is a diffeomorphism from a sufficiently small neighborhood of $0\in T_pM$; it is a local statement only.\\
Geodesic completeness & Every maximal geodesic extends for all affine time; it does not imply that $\exp_p$ is globally injective.\\
Metric completeness & Every metric Cauchy sequence converges; its relation to geodesic completeness is a further theorem with stated Riemannian hypotheses.\\
Minimizing behavior & A radial geodesic minimizes on a sufficiently short range; cut and conjugate phenomena can end global minimizing behavior.\\
\bottomrule
\end{tabular}
\caption{The local exponential-map theorem, geodesic completeness, metric completeness, and minimizing behavior are distinct assertions and require separate hypotheses.}
\label{tab:vii23-c15-completeness-distinction}
\end{table}
